
\chapter{Introduzione}~\label{introduzione}
................

\section{Comandi vari}\label{comandi}


La sezione si richiama all'interno del testo con: Sezione~\ref{comandi}


\begin{itemize}
\item lista con *
\item ....
\end{itemize}


\begin{enumerate}
\item lista numerata
\item ....
\end{enumerate}

\begin{definition}
mettere una definizione
\end{definition}

\begin{teo}
mettere un teorema
\end{teo}
\pf la prova
\ep


Per fare tabelle e richiamarle Tabella~\ref{ds1}: 


\begin{table}
\baselineskip .6cm

\hrulefill

{\small

\begin{minipage}[c]{0.3\linewidth}

\begin{tabbing}
\\
scritta \\
scritta \\
\end{tabbing}
\end{minipage}

\hrulefill
\caption{ quello che si vuole scrivere.} \label{ds1}
}
\end{table}



per richiamare una figura nel testo ......Figure~\ref{etichettafigura}


\begin{figure}
 \centering \epsfig{file=nomefigura.eps, width=1.3in}
 \caption{scritta sotto la figura.\label{etichettafigura}}
 \end{figure}

Per allargare o restringere la figura cambiare width.

per creare la figura usare xfig \newline



{\bf scritta in neretto}, {\em scritta in italico}\newline


% commenta la riga \newline

Formule:


$1\leq i\geq l-1$  1<= i >= l-1



